\section{What one finite-field detector would have to certify}
\label{sec:one-prime}
This section distinguishes flatness over the arithmetic base from the relative flatness that would make one finite-field cycle propagate through a whole family.

Let $\mathcal S$ be an integral scheme of finite type, flat over a localization of the ring of integers of a number field. Let $\mathcal A\to\mathcal S$ be a polarized abelian scheme with a fixed integral quadratic action whose characteristic-zero family is of Weil type. Fix an auxiliary prime $\ell$ invertible on $\mathcal S$. The integral polynomial projector of Section~\ref{sec:weil} defines the Weil summand in the corresponding lisse $\Q_\ell$-sheaf. For $D\ge1$, let $H_D\to\mathcal S$ be a projective parameter space for effective middle-dimensional cycles of degree at most $D$. Such a space is obtained by the construction of Lemma~\ref{lem:parameters}, using single cycles instead of pairs.

\begin{proposition}\label{prop:flat-detector-witness}
Suppose a geometric point $h\in H_D$ parametrizes a cycle with nonzero Weil projection, and the morphism $H_D\to\mathcal S$ is flat at $h$. Then every complex fibre of the family has a detector of degree at most $D$. The point $h$ may have positive characteristic.
\end{proposition}

\begin{proof}
Applying the projector to the universal cycle class gives a lisse section on $H_D$. Its nonzero locus $H_D^{W}$ is a union of geometric connected components. It is open and closed, descends over the base by Galois invariance, and is proper over $\mathcal S$. A flat morphism of finite presentation is open near a point where it is flat: restrict to the open flat locus \cite[Theorems~2.4.6 and 11.1.1]{EGAIV}. Thus the image of $H_D^{W}$ contains a nonempty open subset of $\mathcal S$. Properness makes the image closed, and integrality of $\mathcal S$ makes that image all of $\mathcal S$. In each complex fibre, nonzero projection of an effective cycle implies nonzero projection of at least one of its irreducible components, of degree no larger than $D$.
\end{proof}

For example, if a local-complete-intersection detector $Y$ on a geometric finite-field fibre has $H^1(Y,N_{Y/\mathcal A_s})=0$, the usual embedded deformation theory makes the relative Hilbert scheme smooth at $[Y]$. This gives the proposition's flatness hypothesis. The vanishing is sufficient, not necessary; relative flatness can hold even when this obstruction space is nonzero. No such detector with the required relative deformation property is constructed here in an unresolved eightfold family.

Flatness only over the ring of integers does not suffice. Consider the proper closed immersion
\[
 H=\{x=0\}\longrightarrow\mathcal S=\mathbf A^1_\Z.
\]
The scheme $H\cong\Spec\Z$ is flat over $\Z$ and has no vertical components. For any rational prime $p$, the section $s_p:\Spec\Z\to\mathcal S$ defined by $x=p$ has fibre product
\[
 H\times_{\mathcal S,s_p}\Spec\Z\cong\Spec\mathbf F_p.
\]
It has a point at $p$ and no characteristic-zero point. Since $p$ can be arbitrarily large, discarding a fixed set of vertical primes of the universal parameter space does not solve this problem. The obstruction is failure of flatness over the moduli base along the chosen section.

For a fixed CM fibre, Proposition~\ref{prop:lift} still gives a finite exceptional set depending on that fibre and on $D$. A degree-$D$ detector at a place outside that set suffices for that fibre. Proposition~\ref{prop:flat-detector-witness} is a different sufficient condition: a single detector together with relative flatness makes an entire parameter component dominate the moduli base. A construction must establish one of these properties; flatness over the arithmetic base alone establishes neither.
